\chapter{回归分析}
\begin{theorem}[5.1.1]
若对任意 $x$，
\[
E(Y \mid X=x)=a+bx,
\tag{5.1.4}
\]
则
\[
E(Y \mid X=x)
=
\mu_2
+
\rho
\sqrt{\frac{\operatorname{Var}(Y)}{\operatorname{Var}(X)}}
(x-\mu_1),
\]
其中
\[
\mu_1=E(X),\qquad
\mu_2=E(Y),\qquad
\rho=
\frac{\operatorname{Cov}(X,Y)}
{\sqrt{\operatorname{Var}(X)\operatorname{Var}(Y)}}.
\]
\end{theorem}

\section{简单线性回归模型}

\begin{definition}[简单线性回归模型]
设
\[
Y_i = \beta_0 + \beta_1 x_i + \varepsilon_i, \quad i=1,\ldots,n,
\]
其中 $\varepsilon_i \overset{\text{i.i.d.}}{\sim} N(0,\sigma^2)$，$x_i$ 为已知常数，$\beta_0,\beta_1,\sigma^2$ 未知。
\end{definition}

\begin{theorem}[OLS 估计量]
最小二乘估计量为
\[
\hat{\beta}_1 = \frac{\sum_{i=1}^n (x_i-\bar{x})(Y_i-\bar{Y})}{\sum_{i=1}^n (x_i-\bar{x})^2}, \qquad
\hat{\beta}_0 = \bar{Y} - \hat{\beta}_1\bar{x}.
\]
\end{theorem}

\begin{theorem}[点预测]
给定新观测点 $x^*$，因变量的点预测值为
\[
\hat{Y}(x^*) = \hat{\beta}_0 + \hat{\beta}_1 x^*.
\]
\end{theorem}
\begin{theorem}[OLS 估计量]
最小化残差平方和 $Q(\beta_0,\beta_1)=\sum_{i=1}^n(Y_i-\beta_0-\beta_1 x_i)^2$，
对 $\beta_0,\beta_1$ 各偏导令为零：
\[
\frac{\partial Q}{\partial\beta_0}=0
\;\Rightarrow\;\sum(Y_i-\beta_0-\beta_1 x_i)=0,\qquad
\frac{\partial Q}{\partial\beta_1}=0
\;\Rightarrow\;\sum x_i(Y_i-\beta_0-\beta_1 x_i)=0.
\]
解正规方程组得
\[
\hat{\beta}_1=\frac{S_{xy}}{S_{xx}},\qquad
\hat{\beta}_0=\bar{Y}-\hat{\beta}_1\bar{x},
\]
其中 $S_{xx}=\sum(x_i-\bar{x})^2$，$S_{xy}=\sum(x_i-\bar{x})(Y_i-\bar{Y})$。
\end{theorem}

\begin{note}
用平方而非绝对值的理由：$|e_i|$ 不可微，无法求导；
$e_i^2$ 给出闭合解。
更深层：正态误差下 OLS $=$ MLE；Gauss--Markov 定理保证 OLS 是 BLUE。
\end{note}

\begin{dxtips}
斜率 $=$ 协方差 $\div$ 方差：$\hat\beta_1=S_{xy}/S_{xx}=\widehat{\mathrm{Cov}}(x,Y)/\widehat{\mathrm{Var}}(x)$。
截距由"直线过中心点 $(\bar x,\bar Y)$"反推。
\end{dxtips}
\section{投影与幂等矩阵}

\begin{definition}[幂等矩阵]
矩阵 $P$ 满足 $P^2 = P$，则称 $P$ 为幂等矩阵。
\end{definition}

\begin{theorem}[幂等矩阵的二次型服从 $\chi^2$ 分布]
设 $\boldsymbol{X} \sim N(\boldsymbol{0}, \sigma^2 I_n)$，$P$ 为 $n\times n$ 对称幂等矩阵，$\mathrm{rank}(P)=r$，则
\[
\frac{\boldsymbol{X}^T P \boldsymbol{X}}{\sigma^2} \sim \chi^2(r).
\]
\end{theorem}

\begin{example}[投影定理：5.2.2]
设 $\boldsymbol{u} = \dfrac{1}{\sqrt{n}}(1,1,\ldots,1)^T$，则 $\boldsymbol{u}^T\boldsymbol{u}=1$，令 $P = I_n - \boldsymbol{u}\boldsymbol{u}^T$。

\textbf{验证幂等：}
\[
P^2 = (I_n-\boldsymbol{u}\boldsymbol{u}^T)^2 = I_n - 2\boldsymbol{u}\boldsymbol{u}^T + \boldsymbol{u}\underbrace{(\boldsymbol{u}^T\boldsymbol{u})}_{=1}\boldsymbol{u}^T = I_n - \boldsymbol{u}\boldsymbol{u}^T = P. \quad \checkmark
\]

\textbf{验证秩：}$\mathrm{rank}(\boldsymbol{u}\boldsymbol{u}^T)=1$，故 $\mathrm{rank}(P)=n-1$。

由幂等矩阵定理：
\[
\frac{\boldsymbol{X}^T(I_n-\boldsymbol{u}\boldsymbol{u}^T)\boldsymbol{X}}{\sigma^2} \sim \chi^2(n-1).
\]
\end{example}